\section{The operational question}
\label{sec:operational}

Can every $s$ be given an ordered graph that realizes the Tchoukaillon
schedule, with no condition on the paths? This is exercise~16 under the
operational reading, and it is open. This section collects what we know:
explicit examples for small $s$, an example showing that the operational
and basic-row readings differ, structural consequences of the
schedule's first rounds, and two obstructions that come from the input
order.

\subsection{Small sizes}
\label{sec:catalog}

An earlier computer search found ordered graphs realizing the schedule for
every even $s\le40$ (Appendix~\ref{app:catalog}). It extended one fixed
ordered graph step by step, rerunning the whole execution after each
proposed extension and rejecting any extension that changed an earlier
round. Its program is not among the accompanying programs, but the inputs
it found are. The program \texttt{search.py} (Appendix~\ref{app:code}) is a
new search of the same kind. It finds other graphs, for every even
$s\le46$ with its default seed and, with other seeds, also for $s=48$, 52,
54, and~56. With Proposition~\ref{prop:pad} for $s=50$ and the next
proposition for odd~$s$, the schedule can therefore be realized for every
$s\le57$. It starts from the normal form of Appendix~\ref{app:catalog} and
designs one round at a time, in the manner of Knuth's matrices
(answer~15(d)): at level~$L$ it splits the rows into groups of at most
$L+1$, and for each group it adds the edges of one SAP that merges the
group's rows, placed in the lists so that Algorithm~H should find exactly
that SAP. After each proposed group it reruns the whole execution
from~\Hs1, and it rejects the group if an earlier round changes or the
planned SAPs are not found. Neither search is guaranteed to succeed
(with its default seed \texttt{search.py} finds nothing for $s=48$, 50,
or~52 in two minutes each, and five other seeds found nothing for $s=58$, 60,
or~62 in fifteen minutes each), and neither proves anything about
larger~$s$. We have replayed every
example from both searches from \Hs1 with two independent implementations
of Algorithm~H and checked its final matching with a separate
maximum-matching program (Appendix~\ref{app:code}). A max-flow computation
of $\nu$ in every round shows that all these examples are in fact basic-row
realizations.

\begin{proposition}
\label{prop:odd}
A realization for $s=2h$ yields one for $s=2h+1$: add an isolated
girl--boy edge, numbered after all other vertices. The new realization is
basic-row if the old one is.
\end{proposition}

\begin{proof}
\Hs1 matches the new edge first, no search ever reaches it, and
$\floor{2h/2}=\floor{(2h+1)/2}$, so the schedule is unchanged.
\end{proof}

\subsection{Operational is weaker than basic-row}
\label{sec:no-cover}

The matrices of answer~15(d) keep $\nu=f$ in every round
(property~$\Bp$). The schedule does not require this. Take $s=8$ and the
ordered lists
\begin{equation}
  [5,7];\ [3,5];\ [3];\ [1];\ [1,2];\ [3,4];\ [5,2,6];\ [7,8].
  \label{eq:tree8}
\end{equation}
The graph is a tree with 15 edges (Fig.~\ref{fig:tree8}). Step~\Hs1 leaves
girls~$4,3,2,1$ free, and the run is
\begin{center}\small
\begin{tabular}{crrrrr}
\toprule
$L$&$f$&SAPs&$W_5$&$W_{13}$&$W_{11}$\\
\midrule
2&4&2&15&7&8\\
3&2&1&11&3&5\\
4&1&1&7&4&4\\
\bottomrule
\end{tabular}
\end{center}
which is the schedule $4\to2\to1\to0$. After the round at level~2 (round~2
of the run), the free girls are $g_3,g_4$ and the free boys~$b_6,b_8$, and
there are exactly three augmenting paths,
\begin{align*}
 &g_4\sg b_1\db g_5\sg b_2\db g_7\sg b_6 &&(\text{level }3),\\
 &g_3\sg b_3\db g_2\sg b_5\db g_1\sg b_7\db g_8\sg b_8 &&(\text{level }4),\\
 &g_4\sg b_1\db g_5\sg b_2\db g_7\sg b_5\db g_1\sg b_7\db g_8\sg b_8 &&(\text{level }5).
\end{align*}
(In a tree there is at most one path between two vertices, so one only
has to check the four pairs of a free girl and a free boy.) So with two
free girls there is only one SAP, $\nu=1<f$, and property~$\Bp$ fails for
every choice of perfect matching. An operational realization need not be
basic-row, and a search restricted to Knuth-style matrices could miss
solutions.

\begin{figure}[ht]
\centering
\begin{tikzpicture}[]
\begin{scope}[on background layer]
\draw[sap] (6.120,0.000) -- (5.100,0.000) -- (4.080,0.000) -- (3.060,0.000);
\draw[sap] (5.100,-0.902) -- (5.100,-1.804) -- (6.120,-1.804) -- (7.140,-1.804);
\end{scope}
\draw[edge,draw=black!55,line width=.7pt] (6.120,0.000) -- (5.100,0.000);
\draw[edge,draw=black!55,line width=.7pt] (6.120,0.000) -- (7.140,0.000);
\draw[edge,draw=black!55,line width=.7pt] (5.100,-0.902) -- (5.100,-1.804);
\draw[edge,draw=black!55,line width=.7pt] (5.100,-0.902) -- (5.100,0.000);
\draw[edge,draw=black!55,line width=.7pt] (4.080,-1.804) -- (5.100,-1.804);
\draw[edge,draw=black!55,line width=.7pt] (0.000,0.000) -- (1.020,0.000);
\draw[couple] (2.040,0.000) -- (1.020,0.000);
\draw[edge,draw=black!55,line width=.7pt] (2.040,0.000) -- (3.060,0.000);
\draw[couple] (6.120,-1.804) -- (5.100,-1.804);
\draw[edge,draw=black!55,line width=.7pt] (6.120,-1.804) -- (7.140,-1.804);
\draw[couple] (4.080,0.000) -- (5.100,0.000);
\draw[edge,draw=black!55,line width=.7pt] (4.080,0.000) -- (3.060,0.000);
\draw[edge,draw=black!55,line width=.7pt] (4.080,0.000) -- (4.080,-0.902);
\draw[couple] (8.160,0.000) -- (7.140,0.000);
\draw[edge,draw=black!55,line width=.7pt] (8.160,0.000) -- (9.180,0.000);
\node[girl,free] at (0.000,0.000) {4};
\node[girl] at (2.040,0.000) {5};
\node[girl] at (4.080,0.000) {7};
\node[girl,free] at (6.120,0.000) {1};
\node[girl] at (8.160,0.000) {8};
\node[girl,free] at (5.100,-0.902) {2};
\node[girl,free] at (4.080,-1.804) {3};
\node[girl] at (6.120,-1.804) {6};
\node[boy] at (1.020,0.000) {1};
\node[boy,free] at (3.060,0.000) {2};
\node[boy] at (5.100,0.000) {5};
\node[boy] at (7.140,0.000) {7};
\node[boy,free] at (9.180,0.000) {8};
\node[boy,free] at (4.080,-0.902) {6};
\node[boy] at (5.100,-1.804) {3};
\node[boy,free] at (7.140,-1.804) {4};
\node[font=\footnotesize] at (4.6,-2.45) {after \Hs1: $f=4$, $L=2$};
\end{tikzpicture}

\bigskip

\begin{tikzpicture}[]
\begin{scope}[on background layer]
\draw[sap] (0.000,0.000) -- (1.020,0.000) -- (2.040,0.000) -- (3.060,0.000) -- (4.080,0.000) -- (4.080,-0.902);
\draw[sap,draw=boydraw,opacity=.2] (4.080,-1.804) -- (5.100,-1.804) -- (5.100,-0.902) -- (5.100,0.000) -- (6.120,0.000) -- (7.140,0.000) -- (8.160,0.000) -- (9.180,0.000);
\end{scope}
\draw[couple] (6.120,0.000) -- (5.100,0.000);
\draw[edge,draw=black!55,line width=.7pt] (6.120,0.000) -- (7.140,0.000);
\draw[couple] (5.100,-0.902) -- (5.100,-1.804);
\draw[edge,draw=black!55,line width=.7pt] (5.100,-0.902) -- (5.100,0.000);
\draw[edge,draw=black!55,line width=.7pt] (4.080,-1.804) -- (5.100,-1.804);
\draw[edge,draw=black!55,line width=.7pt] (0.000,0.000) -- (1.020,0.000);
\draw[couple] (2.040,0.000) -- (1.020,0.000);
\draw[edge,draw=black!55,line width=.7pt] (2.040,0.000) -- (3.060,0.000);
\draw[edge,draw=black!55,line width=.7pt] (6.120,-1.804) -- (5.100,-1.804);
\draw[couple] (6.120,-1.804) -- (7.140,-1.804);
\draw[edge,draw=black!55,line width=.7pt] (4.080,0.000) -- (5.100,0.000);
\draw[couple] (4.080,0.000) -- (3.060,0.000);
\draw[edge,draw=black!55,line width=.7pt] (4.080,0.000) -- (4.080,-0.902);
\draw[couple] (8.160,0.000) -- (7.140,0.000);
\draw[edge,draw=black!55,line width=.7pt] (8.160,0.000) -- (9.180,0.000);
\node[girl,free] at (0.000,0.000) {4};
\node[girl] at (2.040,0.000) {5};
\node[girl] at (4.080,0.000) {7};
\node[girl] at (6.120,0.000) {1};
\node[girl] at (8.160,0.000) {8};
\node[girl] at (5.100,-0.902) {2};
\node[girl,free] at (4.080,-1.804) {3};
\node[girl] at (6.120,-1.804) {6};
\node[boy] at (1.020,0.000) {1};
\node[boy] at (3.060,0.000) {2};
\node[boy] at (5.100,0.000) {5};
\node[boy] at (7.140,0.000) {7};
\node[boy,free] at (9.180,0.000) {8};
\node[boy,free] at (4.080,-0.902) {6};
\node[boy] at (5.100,-1.804) {3};
\node[boy] at (7.140,-1.804) {4};
\node[font=\footnotesize] at (4.6,-2.45) {after round 2: $f=2$, $L=3$};
\end{tikzpicture}

\caption{An operational realization for $s=8$ that is not basic-row,
labelled by the numbers of~\eqref{eq:tree8}. Top: the matching after~\Hs1
and the two SAPs of the round at level~2. Bottom: the next matching. Its
only SAP is shaded orange; the faint blue path is the augmenting path of
level~4. (The third augmenting path, of level~5, is not drawn.)}
\label{fig:tree8}
\end{figure}

\subsection{The shape forced by the first matching}
\label{sec:blocks}

For even $s=2h$ the schedule asks~\Hs1 to leave $h$ girls free, so~\Hs1 must
produce a maximal matching with only $h$ couples, the fewest possible. The
next proposition describes what that implies for the graph, whatever the
input order. Write $G^\circ$ and $G^\bullet$ for the girls who are free and
matched after~\Hs1, and $B^\bullet$ and $B^\circ$ for the boys who are
matched and free after~\Hs1.

\begin{proposition}
\label{prop:blocks}
Suppose $s=M=N=2h$ and~\Hs1 leaves $h$ girls free. Then:
\begin{enumerate}[label=(\alph*)]
\item there are no edges between $G^\circ$ and $B^\circ$, so the biadjacency
  matrix has the block form of Fig.~\ref{fig:blocks};
\item every perfect matching pairs $G^\circ$ with $B^\bullet$ and $G^\bullet$
  with $B^\circ$;
\item whenever $f$ girls are free later on, exactly $f$ couples join
  $G^\bullet$ to $B^\bullet$, and $h-f$ couples are of each of the types
  $G^\circ$--$B^\bullet$ and $G^\bullet$--$B^\circ$;
\item relative to any perfect matching, every basic path contains exactly
  one couple of type $G^\bullet$--$B^\bullet$.
\end{enumerate}
\end{proposition}

\begin{proof}
(a) is maximality of the matching from~\Hs1. For (b), the $h$ girls of
$G^\circ$ can only be matched into the $h$ boys of $B^\bullet$, which leaves
$B^\circ$ for $G^\bullet$. (c) Vertices never become free again once
matched, so $G^\bullet$ and $B^\bullet$ stay covered. If $x,y,z$ count the
current couples of types $G^\circ$--$B^\bullet$, $G^\bullet$--$B^\bullet$,
$G^\bullet$--$B^\circ$, then $x+y=|B^\bullet|=h$, $y+z=|G^\bullet|=h$, and
$x+y+z=2h-f$; hence $y=f$. (d) A basic path starts at a free girl in
$G^\circ$ and ends at a free boy in $B^\circ$. Its perfect-matching edges
stay within $G^\circ\cup B^\bullet$ or within $G^\bullet\cup B^\circ$
by~(b), its couples of the two mixed types do too, and a couple of type
$G^\bullet$--$B^\bullet$ is the only way to pass from the first part to the
second; by~(a) there is no way back.
\end{proof}

\begin{figure}[ht]
\centering
% Adjacency blocks after a minimum H1 (rows: girls; columns: boys).
\begin{tikzpicture}[x=1.15cm,y=1.0cm]
\draw[fill=boyfill] (0,0) rectangle (1,1);
\draw[fill=white] (1,0) rectangle (2,1);
\draw[fill=boyfill] (0,-1) rectangle (1,0);
\draw[fill=boyfill] (1,-1) rectangle (2,0);
\node at (1.5,.5) {$0$};
\node[font=\small,anchor=east] at (-.1,.5) {$G^\circ$ (free)};
\node[font=\small,anchor=east] at (-.1,-.5) {$G^\bullet$ (matched)};
\node[font=\small,anchor=south] at (.5,1.05) {$B^\bullet$};
\node[font=\small,anchor=south] at (1.5,1.05) {$B^\circ$};
\node[font=\footnotesize,anchor=south,text=black!65] at (.5,1.45) {matched};
\node[font=\footnotesize,anchor=south,text=black!65] at (1.5,1.45) {free};
\end{tikzpicture}

\caption{After a minimum-size~\Hs1 matching no free girl is adjacent to a
free boy: the biadjacency matrix has a zero block.}
\label{fig:blocks}
\end{figure}

In terms of Knuth's matrices, each row consists of some indices of
initially free girls followed by some indices of initially matched girls.
For the 16-pair example this is the remark at the end of answer~15(d):
an edge $b_{\rm even}\sg g_{\rm odd}$ would join a free girl to a free boy
of the first matching.

\subsection{Realizations need cycles}

The trees of Section~\ref{sec:sparse} already force Algorithm~H to do
$\Theta(E\sqrt V)$ work. No forest, however, can realize the schedule for
large $s$: its first two rounds force many independent cycles. The
argument uses only the prescribed numbers of SAPs in those rounds, so it
applies to any algorithm that removes a maximal set of disjoint SAPs per
round, in any order.

\begin{lthm}{C}
Let an ordered graph with $s=M=N=2h>0$ realize the Tchoukaillon schedule.
Let $a_2=\ceil{h/3}$ and $a_3=\ceil{(h-a_2)/4}$ be the numbers of SAPs in
its rounds at levels~2 and~3, and let $p=3a_2-h\in\{0,1,2\}$. Then
\[
   E\ \ge\ 3h+4a_2-8p-a_3 .
\]
Consequently the cycle rank $\beta=E-V+\kappa$, where $\kappa$ is the number
of connected components, satisfies $\beta\ge a_2-a_3-7p+1\ge\ceil{(h-76)/6}$.
In particular the graph is not a forest if $h\ge77$, or if $3\mid h$; and for
large $s$,
\[
   E\ge\tfrac{25}{12}\,s-O(1),\qquad \beta\ge\tfrac1{12}\,s-O(1).
\]
\end{lthm}

\begin{proof}
Fix a perfect matching $P$ and let $M_1$ be the matching produced by~\Hs1.

\emph{Basic paths after~\Hs1.} The $2h$ endpoints of $M_1$ must touch all
$2h$ edges of $P$, by maximality, so each edge of $P$ has exactly one
matched endpoint. Hence $P\cup M_1$ is a disjoint union of $h$ basic paths
of level~2, each of the form $u\sg v\db w\sg z$ with $u\in G^\circ$, $v\in
B^\bullet$, $w\in G^\bullet$, $z\in B^\circ$; together they cover all
vertices and use $3h$ edges.

\emph{Good triples.} A SAP of the first dag round also has the form
$u\sg v\db w\sg z$, and it meets at most three basic paths: the one
containing~$u$, the one containing the couple $v\db w$, and the one
containing~$z$. The round selects $a_2$ SAPs, and by maximality they meet
all $h$ basic paths. Call a selected SAP \emph{good} if it meets three basic
paths that no other selected SAP meets. If $q_1,\ldots,q_{a_2}$ are the
numbers of basic paths met by the selected SAPs, put
$\delta=3a_2-\sum q_i\ge0$ and $\omega=\sum q_i-h\ge0$, so that
$\delta+\omega=p$. At most $\delta$ selected SAPs meet fewer than three
paths; and a basic path met by $k\ge2$ selected SAPs contributes $k-1$ to
$\omega$ and involves at most $k\le2(k-1)$ of them. So the number $g$ of
good SAPs satisfies $g\ge a_2-2p$.

\emph{What a good triple becomes.} Number the three basic paths met by a good
SAP as $u_k\sg v_k\db w_k\sg z_k$ ($k=1,2,3$), so that the SAP is
$u_1\sg v_2\db w_2\sg z_3$. After augmenting, the six pairs of the triple
form two basic paths of level~3 (Fig.~\ref{fig:goodtriple}):
\[
  u_2\sg v_2\db u_1\sg v_1\db w_1\sg z_1\qquad\text{and}\qquad
  u_3\sg v_3\db w_3\sg z_3\db w_2\sg z_2 .
\]
These are SAPs of the round at level~3, and their vertices lie in the dag
at the levels given by their positions (a shorter route to one of them,
followed by the rest of the path, would contain an augmenting path of
level~2, and there is none after the first dag round). The triple contains
$9+2=11$ edges of $P\cup M_1$ and of the selected SAP: the ten edges of the
two new paths, and the old couple $v_2w_2$. That edge joins the girl~$w_2$,
at position~3 of the second path, to the boy~$v_2$, at position~1 of the
first, so it points backwards and is not an arc of the level-3 dag.

\emph{The second round.} Let $F$ be the set of edges outside $P\cup M_1$ and
outside the $2g$ selected edges of good triples; $|F|=E-3h-2g$. In the
level-3 dag with $F$ removed, each new path of a good triple is cut off
from everything else: edges of $P\cup M_1$ never leave their triple, the
selected edges stay inside it, and the only remaining edge between the two
new paths is not an arc. So a SAP that meets $q$ of these $2g$ paths uses at
least $q-1$ edges of~$F$. The round at level~3 selects $a_3$ vertex-disjoint
SAPs, which must meet all $2g$ new paths of good triples. Therefore
$2g-a_3\le|F|=E-3h-2g$, and $E\ge3h+4g-a_3\ge3h+4a_2-8p-a_3$.

(When $h=1$ there is no second round, $a_3=0$, and the bound is trivial.)
The bound on $\beta$ follows from $V=4h$, $\kappa\ge1$, and $h=3a_2-p$; the
explicit form comes from substituting $a_3=\ceil{(2a_2-p)/4}$ and checking
$p=0,1,2$, the weakest case being $p=2$ with $a_2$ even. For $p=0$ one gets
$\beta\ge\floor{a_2/2}+1>0$. The asymptotic statements follow from
$a_2=h/3+O(1)$ and $a_3=h/6+O(1)$.
\end{proof}

\begin{figure}[ht]
\centering
% A good triple (proof of Theorem C): before and after the first dag round.
\begin{tikzpicture}[x=1.05cm,y=.9cm]
% ---------- left: three basic paths after H1, one selected SAP
\foreach \k/\y in {1/0,2/-1,3/-2}{
  \node[girl,free] (u\k) at (0,\y) {$u_\k$};
  \node[boy]       (v\k) at (1,\y) {$v_\k$};
  \node[girl]      (w\k) at (2,\y) {$w_\k$};
  \node[boy,free]  (z\k) at (3,\y) {$z_\k$};
  \draw[pedge] (u\k)--(v\k) (w\k)--(z\k);
  \draw[couple] (v\k)--(w\k);}
\begin{scope}[on background layer]
  \draw[sap] (u1.center)--(v2.center)--(w2.center)--(z3.center);
\end{scope}
\draw[edge,draw=black!60,line width=.7pt] (u1)--(v2) (w2)--(z3);
\node[font=\footnotesize,anchor=north] at (1.5,-2.5) {after \Hs1};
% ---------- right: the two new basic paths
\begin{scope}[xshift=5.3cm]
  \foreach \p in {1,2,3}{\node[font=\footnotesize,text=black!60] at ({2*\p-1.5},.65) {position \p};}
  \node[girl,free] (nu2) at (0,0) {$u_2$};
  \node[boy]       (nv2) at (1,0) {$v_2$};
  \node[girl]      (nu1) at (2,0) {$u_1$};
  \node[boy]       (nv1) at (3,0) {$v_1$};
  \node[girl]      (nw1) at (4,0) {$w_1$};
  \node[boy,free]  (nz1) at (5,0) {$z_1$};
  \node[girl,free] (nu3) at (0,-1.9) {$u_3$};
  \node[boy]       (nv3) at (1,-1.9) {$v_3$};
  \node[girl]      (nw3) at (2,-1.9) {$w_3$};
  \node[boy]       (nz3) at (3,-1.9) {$z_3$};
  \node[girl]      (nw2) at (4,-1.9) {$w_2$};
  \node[boy,free]  (nz2) at (5,-1.9) {$z_2$};
  \draw[pedge] (nu2)--(nv2) (nu1)--(nv1) (nw1)--(nz1) (nu3)--(nv3) (nw3)--(nz3) (nw2)--(nz2);
  \draw[couple] (nv2)--(nu1) (nv1)--(nw1) (nv3)--(nw3) (nz3)--(nw2);
  \draw[edgecol,line width=.8pt,densely dashed] (nw2) to[bend left=8] (nv2);
  \node[font=\footnotesize,text=black!65,fill=white,inner sep=1.5pt] at (2.55,-.95) {old couple $v_2w_2$};
  \node[font=\footnotesize,anchor=north] at (2.5,-2.5) {after the first dag round};
\end{scope}
\draw[-{Stealth},black!40,line width=.8pt] (3.65,-1)--(4.75,-1);
\end{tikzpicture}

\caption{A good triple in the proof of Theorem~C. Single lines are edges of
the perfect matching $P$, double lines are couples. One SAP of the first
dag round (shaded) meets three basic paths; afterwards their six pairs form
two basic paths of level~3. The old couple $v_2w_2$ runs from position~3
back to position~1, so the level-3 dag cannot use it, and any SAP that
visits both new paths must use some other edge.}
\label{fig:goodtriple}
\end{figure}

\begin{remark}
For odd $s=2h+1$ the same method should give
$E\ge3h+4a_2-8p-a_3-6$ and $\beta\ge\ceil{(h-124)/6}$, so that no forest
realizes the schedule when $s\ge251$. In outline: \Hs1 now leaves $h$ girls
free, and either all $h$ basic paths have level~2 and one couple of $M_1$
belongs to $P$, or one basic path has level~3 and the others level~2. We
have not written out the second case in full (one must check that SAPs
touching the long basic path do not spoil the isolation argument), so we
state this only as a remark.
\end{remark}

The case~$p=0$ also shows how property~$\Up$ fails in such a realization.

\begin{corollary}
If a realization has $6\mid s$, then just after its first dag round the
SAPs have level~3 (five edges), and there is also an augmenting path with
at least nine edges.
\end{corollary}

\begin{proof}
With $p=0$ all $a_2$ selected SAPs are good, so no basic path is left
untouched, and the $2a_2$ new basic paths are disjoint SAPs of level~3,
each meeting a selected SAP. Apply Proposition~\ref{prop:longpaths} with
$q=2a_2$ and $a=a_2$.
\end{proof}

\subsection{The input order matters}
\label{sec:order}

One plan for the operational question is to find a graph with the right
sequence of maximal SAP sets, and then look for an input order that makes
Algorithm~H pick those sets. The plan can fail, because Algorithm~H does
not choose its maximal sets freely: which SAP it finds depends on the
order in \A{QUEUE}, on head insertion in step~\Hs6, and on the replacement
rule of step~\Hs{16}.

\paragraph{A six-pair graph.}
Let the girls~$g_1,\ldots,g_6$ have the neighbour sets
\[
  \{b_1\};\quad\{b_1,b_2,b_5\};\quad\{b_3\};\quad\{b_2,b_3,b_4\};\quad\{b_3,b_5\};\quad\{b_5,b_6\}.
\]
The edges $g_kb_k$ form a perfect matching. Starting from the matching
$b_1g_2$, $b_3g_4$, $b_5g_6$, the three augmentations
\begin{gather*}
  g_5\sg b_3\db g_4\sg b_2,\qquad
  g_1\sg b_1\db g_2\sg b_5\db g_6\sg b_6,\\
  g_3\sg b_3\db g_5\sg b_5\db g_2\sg b_2\db g_4\sg b_4
\end{gather*}
are SAPs at levels~2, 3, 4, and after each of them no SAP of the same level
remains. So the schedule for $s=6$, which removes one SAP at each of these
levels, is consistent with this graph.

\begin{proposition}
No numbering and no list order make Algorithm~H follow this schedule on
this graph.
\end{proposition}

\begin{proof}
A maximal matching with three couples must have exactly one endpoint on
each edge $g_kb_k$ and contain none of them; the only such matching is
$b_1g_2$, $b_3g_4$, $b_5g_6$, so~\Hs1 must produce it, with free girls
$g_1,g_3,g_5$. One SAP can be a maximal set only if it meets all three
basic paths, and the only such SAP is $g_5\sg b_3\db g_4\sg b_2$. Before
the first augmentation every root of the search succeeds, so this must be
the SAP found from the first root, $b_2$, by taking the first arc of each
list. Head insertion at $b_3$ then requires $g_3$ to precede $g_5$ in
\A{QUEUE}, and head insertion at $b_2$ requires $g_2$ to precede $g_4$ at
level~1. Girl~$g_2$ is queued when $g_1$ discovers $b_1$, and $g_4$ when the
earlier of $g_3,g_5$ discovers $b_3$; so $g_1$ comes first, and the queue
must be $(g_1,g_3,g_5)$. Step~\Hs{16} then removes the last girl, $g_5$,
leaving $(g_1,g_3)$. In the next dag, level~1 is $(g_2,g_5)$, both adjacent
to $b_5$, so $\A{BLINK}[b_5]=(g_5,g_2)$ whatever the lists. Whichever of the
free boys~$b_4,b_6$ is searched first finds one of the basic paths
$g_1\sg b_1\db g_2\sg b_2\db g_4\sg b_4$ or $g_3\sg b_3\db g_5\sg b_5\db
g_6\sg b_6$; they are disjoint, so the other is found as well, and the round
removes two SAPs instead of one.
\end{proof}

An exhaustive check of all $6!\,(3!)^2(2!)^2=103680$ girl numberings and
list orders (the boy numbering plays no role) agrees: 740 of them give a
three-couple~\Hs1 matching, and none of those follows the schedule. Other
six-pair graphs do realize it (Appendix~\ref{app:catalog}).

\paragraph{Triangular chain graphs.}
Let $\Delta_{2h}$ have girls and boys~$1,\ldots,2h$, with $g_j\sg b_i$ an
edge whenever $i\le j$. Hegerfeld and Kratsch show that chain graphs of
this kind can need many phases when the maximal sets are chosen
adversarially~\cite[Section~5.3]{HegerfeldKratsch}. With Algorithm~H and a
minimum-size first matching, the order of the search rules out the
Tchoukaillon schedule whenever $3\mid h$:

\begin{proposition}
\label{prop:chain}
Give $\Delta_{2h}$ any numbering and any list order. If~\Hs1 leaves $h$ girls
free and the first dag round finds $a$ SAPs, then $3a>h$. So when $3\mid h$,
the graph $\Delta_{2h}$ cannot realize the schedule, whose first round would
remove exactly $h/3$ SAPs.
\end{proposition}

\begin{proof}
As in the previous proof, each edge $g_kb_k$ has exactly one matched
endpoint. Let $T$ be the set of $k$ with $g_k$ free (so $b_k$ is matched)
and $S$ its complement (so $b_k$ is free). Maximality forbids an edge from a
free girl~$g_t$ to a free boy~$b_u$; since $g_tb_u$ is an edge when $u\le
t$, we must have $u>t$ for all $t\in T$ and $u\in S$. Hence
$T=\{1,\ldots,h\}$ and $S=\{h+1,\ldots,2h\}$. Write $u_i=g_i$, $v_i=b_i$,
$x_j=g_{h+j}$, $z_j=b_{h+j}$ for $1\le i,j\le h$; the matching from~\Hs1 is
$x_j\db v_{\pi(j)}$ for a permutation~$\pi$, and the SAPs are exactly the
paths
\[
   u_i\sg v_{\pi(j)}\db x_j\sg z_d\qquad\text{with }\pi(j)\le i\text{ and }d\le j .
\]
Suppose $a<h$ SAPs are selected, let $i_0$ be the largest index of an
unselected free girl and $d_0$ the smallest index of an unselected free boy.
An unselected $x_j$ with $\pi(j)\le i_0$ and $j\ge d_0$ would give a SAP
disjoint from the selected ones; so each of the $h-a$ unselected $x_j$ has
$\pi(j)>i_0$ or $j<d_0$. At most $h-i_0\le a$ of them satisfy the first
condition and at most $d_0-1\le a$ the second; hence $h\le3a$.

Suppose $h=3a$. Then all these inequalities are equalities: the selected
free girls are $u_{2a+1},\ldots,u_{3a}$, the selected free boys are
$z_1,\ldots,z_a$, every selected $x_j$ has $\pi(j)\le2a$ and $j>a$, and
every $x_j$ with $\pi(j)>2a$ is unselected and has $j>a$. Girl~$u_{2a}$ is
never selected. Each selected SAP passes through a boy~$v_{\pi(j)}$ with
$\pi(j)\le2a$, whose list contains $u_{2a}$; that arc cannot have been
consumed earlier, because following it would have produced a SAP ending at
$u_{2a}$. Since lists are in reverse queue order, every selected free girl
comes after $u_{2a}$ in \A{QUEUE}. Only the selected free girls are adjacent
to the boys~$v_{2a+1},\ldots,v_{3a}$, while $u_{2a}$ is adjacent to all of
$v_1,\ldots,v_{2a}$; so the latter are all discovered first, and their
mates come first in level~1. In particular every $x_j$ with $\pi(j)>2a$ is
queued after every selected $x_j$. Now consider the first root of the
search. Before any augmentation every root succeeds, so it is a selected
free boy~$z_d$ with $d\le a$. All $x_j$ with $\pi(j)>2a$ are adjacent to it
(since $j>a\ge d$) and come later in the queue, so the head of
$\A{BLINK}[z_d]$ is not a selected girl; and the search from it succeeds.
So the first SAP found uses an $x_j$ that should have remained unselected,
a contradiction.
\end{proof}

For $h=3$ the schedule is consistent with the graph itself: starting from
the matching $x_1v_1$, $x_2v_3$, $x_3v_2$, the graph $\Delta_6$ admits the
single maximal SAPs
\begin{gather*}
  u_3\sg v_2\db x_3\sg z_1,\qquad
  u_1\sg v_1\db x_1\sg z_1\db x_3\sg z_2,\\
  u_2\sg v_2\db u_3\sg v_3\db x_2\sg z_2\db x_3\sg z_3
\end{gather*}
at levels~$2,3,4$. It is the order of Algorithm~H's search that rules them
out. Other graphs with $s=6$ do realize the schedule.

\subsection{What an induction would need}

One might try to pass from $s=2h$ to $s=2h+2$ by adding two pairs. The
schedule changes in a very controlled way. Write $f_L(h)$ and $a_L(h)$ for
the schedule with $f_2(h)=h$, extended by zeros.

\begin{lemma}
\label{lem:step}
Let $K$ be the first level~$L\ge2$ such that $L+1$ divides $f_L(h)$. Then
$f_L(h+1)=f_L(h)+1$ for $2\le L\le K$ and $f_L(h+1)=f_L(h)$ for $L>K$; so
$a_L(h+1)=a_L(h)+\Iv{L=K}$.
\end{lemma}

\begin{proof}
The map $x\mapsto x-\ceil{x/(L+1)}$ increases by $1$ from $x$ to $x+1$,
except when $(L+1)\mid x$, where it stays the same. Start from a difference
of $1$ at level~2 and follow it.
\end{proof}

So adding two pairs should add one free girl, keep every count of SAPs
except one extra at level~$K$, and then rejoin the old schedule. When $K=2$,
that is, when $3\mid h$, this can be done.

\begin{proposition}
\label{prop:pad}
Let an ordered graph $G$ realize the schedule for $s=2h$, where $3\mid h$.
Let $q$ be the first girl in \A{QUEUE} after step~\Hs1, let $b$ be the first
boy on her list, and let $x$ be the girl whom~\Hs1 matches to~$b$. Form $G^+$
by adding girls $u$ and $w$, numbered below all girls of~$G$, boys $v$
and~$z$, the lists $u\colon[v]$ and $w\colon[v,z]$, and the edge $xz$, placed
in $x$'s list immediately after~$b$. Then $G^+$ realizes the schedule for
$s+2$, and it is basic-row if $G$ is.
\end{proposition}

\begin{proof}
Step~\Hs1 scans the girls of $G$ first and does what it did in~$G$: when $x$
is scanned, the boys before $b$ in her list are taken and $b$ is free, so
she never reaches~$z$. Then $w$ takes~$v$, and $u$ remains free. So
\A{QUEUE} is that of~$G$ followed by~$u$, with $h+1$ girls.

In the first dag, $q$ is scanned first. Her first boy $b$ is matched,
because $q$ is free after a maximal matching; so $b$ is the first boy of
level~1, and $x$ is the first girl of level~1. In her list, $x$ meets
matched boys (all taken by~\Hs1), then her mate~$b$, then $z$, who is free:
so $L=2$ and $z$ is the first root. The other free boys of level~2 are
entered in the same order as in~$G$, because only matched boys, which are
never roots, are treated differently once $L$ is known. Girl~$u$ is last at
level~0, so $w$ is last at level~1, and step~\Hs6 puts her at the head of
$\A{BLINK}[z]$. The search from~$z$ therefore takes~$w$ and then~$v$, whose
only girl at level~0 is~$u$: it finds the SAP $u\sg v\db w\sg z$. Since $u$
is the last free girl, step~\Hs{16} moves nobody. From then on the roots,
the lists \A{BLINK} of the boys of~$G$ (none of which contains a new girl),
and \A{QUEUE} are as in~$G$, so the round finds the $a_2(h)$ SAPs of~$G$ in
the same order, with the same moves in step~\Hs{16}. By
Lemma~\ref{lem:step} with $K=2$, that makes $a_2(h+1)$ SAPs, and after this
round the two schedules agree.

In later rounds $u$, $v$, $w$, $z$ are matched, and the only edge between them
and~$G$ is~$xz$. The breadth-first search may enter $z$, $w$, $v$, $u$
from~$x$, but they lead to no free boy, they do not change the relative
order of the vertices of~$G$ in \A{MARKED} or \A{QUEUE}, and no list
$\A{BLINK}[b']$ with $b'$ in~$G$ gains a girl. The depth-first search starts
at free boys of~$G$, and from a boy of~$G$ it reaches only girls of~$G$,
whose mates are boys of~$G$; so it never leaves~$G$. Every later round is
therefore the corresponding round of~$G$. Finally, in the first dag round
the $h$ disjoint SAPs of a basic-row~$G$ and the path $u\sg v\db w\sg z$
give $\nu=h+1=f$, and afterwards no augmenting path meets the new vertices.
\end{proof}

The edge $xz$ matters. Without it $z$ is the last root, the new SAP is found
last, and by then step~\Hs{16} may have moved~$u$ into the place of a girl
of~$G$, which changes the order of later searches; for $s=42$ and $s=54$
the padding then fails. For $K>2$ the extra free girl must survive several
rounds, and doing this with a fixed set of edges, while preserving every
earlier distance and every earlier \A{BLINK} order, is what we have not
managed: an edge added for a later round can create a shorter path in an
earlier one, or change the order in which an earlier search discovers
vertices.
